% =====================================================================
\section{A Calculus of Guardrail Hierarchies}
\label{ch:calculus}
% =====================================================================

\begin{bigidea}
To reason about stacks, we need to talk about \emph{any} stack and \emph{any} behavior of its parts. So we invent a tiny language with just four building blocks (\textbf{deny}, \textbf{gate}, \textbf{guard}, \textbf{switch}), and a precise definition of what ``run this stack on one proposal'' and ``every stage keeps its contract'' mean. Then \textbf{safe} becomes a mathematical statement: \emph{for every contract-respecting behavior, an accepted proposal satisfies all the goals.}
\end{bigidea}

\subsection{Why a ``calculus''?}

A \emph{calculus} here means: a small formal language with a precise meaning, which strips a problem to its essentials so that it can be reasoned about with proofs. The lambda calculus does this for functions. The paper's calculus does this for stacks of safety checks.

We want to throw away everything about how a classifier or an LLM is implemented. We keep only three things: (1) what a stage \emph{can do to the decision} (stop, accept, pass on, route), (2) what it \emph{assumes}, and (3) what it \emph{promises}. If a theorem holds in this stripped-down world, it holds for every implementation of the stages, including adversarial ones.

\subsection{The ingredients}

\begin{defbox}{goals}
A finite set $\Goals$ of \textbf{goals}: the properties every \emph{accepted} proposal must have. Example: $\Goals=\{\varphi_1,\varphi_2,\varphi_{\mathrm{env}},\varphi_3,\varphi_4\}$ for DevAgent.
\end{defbox}

\begin{defbox}{obligations}
A set of names called \textbf{obligations}: assumptions that a stage's contract may depend on. Examples: \ob{miss\_rate\_bounded}, \ob{safe\_on\_this\_text}, \ob{human\_attentive}. They are just names; what each means is given in English next to the specification, and someone has to be responsible for making it true.
\end{defbox}

\begin{defbox}{layer}
A \textbf{layer} $\ell$ is the \emph{interface} of a learned or human component. It declares two sets:
\begin{itemize}
\item $\req(\ell)\subseteq$ obligations: what its contract \emph{assumes};
\item $\scope(\ell)\subseteq\Goals$: the goals it \emph{covers} when it accepts.
\end{itemize}
There is also a special \textbf{empty layer} $\epsilon$ with $\req(\epsilon)=\emptyset$ and $\scope(\epsilon)=\emptyset$, used for deterministic checks.
\end{defbox}

Notice that we do not say how a layer decides. That is intentional.

\subsection{The four building blocks}

The calculus represents a stage's effect on a proposal using four building blocks:

\begin{enumerate}
\item The proposal can be \textbf{denied} (stopped for good). This is the end of any failed path.
\item A stage can \textbf{check something and stop the proposal if the check fails}, but never accept it by itself. (A rule, or a classifier used as a veto.) We call it a \textbf{gate}.
\item A stage can \textbf{make the final decision}: accept, reject, or pass it to a later stage (escalate). We call it a \textbf{guard}.
\item The system can \textbf{route} the proposal one way or the other depending on a condition (``is the input tainted?''). We call it a \textbf{switch}.
\end{enumerate}

Formally, a hierarchy $H$ is a term of the grammar
\begin{align*}
H ::={}& \deny \mid \gate(n,\varphi,\ell,\Omega,D,H)\\
\mid{}& \guard(n,\ell,\Omega,H) \mid \switch(n,H,H)
\end{align*}
Read ``$::=$'' as ``is one of'', and ``$\mid$'' as ``or''. Let us decode each piece.

\begin{center}
\renewcommand{\arraystretch}{1.3}
\begin{tabularx}{\textwidth}{@{}p{4.8cm}X@{}}
\toprule
\textbf{Piece} & \textbf{Meaning}\\
\midrule
$n$ & A \textbf{node identifier}, the stage's name (e.g.\ $g_1$, $n_3$). Used so that a model can describe each stage's behavior separately.\\
$\ell$ & The stage's \textbf{layer interface}: $\req(\ell)$ gives its assumptions and $\scope(\ell)$ its covered goals.\\
$\Omega$ & The \textbf{evidence}: the set of obligations that the \emph{specification declares} to hold for this one invocation. A \emph{declaration}, not a fact: nothing in the term makes it true.\\
$D$ (gates only) & The \textbf{facts a gate provides}: if the gate passes and its contract applies, these obligations hold at later gate and guard invocations on this run.\\
$\varphi$ (gates only) & The goal the gate \textbf{establishes} if it passes and its contract applies (may be absent, written $\bot$).\\
the last argument $H$ & The \textbf{continuation}: what comes next.\\
\bottomrule
\end{tabularx}
\end{center}

\subsubsection{Gate: can only stop}

A gate can only \emph{stop} a proposal. If the proposal passes the gate, \emph{and the gate's contract applies}, then the gate establishes the goal $\varphi$ and provides the facts $D$. ``The contract applies'' means that all of $\req(\ell)$ is available.

\begin{itemize}
\item If $\ell=\epsilon$ the gate is a \textbf{deterministic gate}: a reference monitor (a small, always-on checker of every action) that decides $\varphi$ \emph{exactly}. Example: ``a secret-labeled file must not go to an external host''.
\item Otherwise it is a \textbf{veto stage}: a learned component, such as a classifier, used only to reject.
\end{itemize}

\subsubsection{Guard: can end the decision}

A guard has a verdict of three kinds: $\acc$ (accept: the decision ends, the proposal is executed), $\blk$ (block: denied), or $\esc$ (escalate: pass to the continuation, typically a costlier stage). \textbf{Every acceptance is the act of a named guard.} We call each guard a \textbf{route}: it is a way a proposal can get through. A stack with no semantic final stage is modelled by a final guard with empty scope that simply releases what reaches it.

\begin{warning}{A guard provides nothing}
When a guard accepts or escalates, \emph{no fact is added} for later stages. A guard's verdict is not evidence. This is the formal version of ``the manager does not learn anything about the lead-lined suitcase from the first operator saying `OK'.'' Only a \emph{gate} whose contract applied can provide facts, and a deterministic gate always applies.
\end{warning}

\subsubsection{Switch: routes, teaches nothing}

A switch chooses between two sub-hierarchies according to a condition (e.g.\ ``is the input tainted?''). Taking a branch establishes no fact, \emph{even if the condition looks like a test}, because a switch has no contract. To learn something from a condition, one must put a \emph{gate} after it.

\subsubsection{Deny: the only leaf}

Leaving the hierarchy denies. A guard that escalates past the end of the hierarchy has nothing to hand the proposal to, so the proposal is denied. This design (``fail closed'') makes sure that the only way to be accepted is for a guard to say $\acc$.

\subsubsection{Build and run a small code-review hierarchy}
\label{sec:code-review-constructors}

\begin{example}{a secret scanner, a router, and two reviewers}
Suppose accepted code changes must satisfy two goals:
\[
\Goals=\{s,r\},\quad s=\text{no exposed credentials},\quad
r=\text{passes the required security review}.
\]
Let $t=\ob{trusted\_input}$ and $q=\ob{secrets\_scanned}$ be obligations. Define interfaces by
\[
\begin{array}{c|cc}
 & \req & \scope\\\hline
L_s & \{t\} & \emptyset\\
L_a & \{q\} & \{r\}\\
L_h & \{q,\ob{human\_attentive}\} & \{r\}
\end{array}
\]
Here $L_s$ is a scanner used as a gate; its goal $s$ and provided fact $q$ are written in the gate itself. Set
\begin{align*}
K_h &=\guard(h,L_h,\{\ob{human\_attentive}\},\deny),\\
K_a &=\guard(a,L_a,\emptyset,\deny),\\
H_{\mathrm{review}}
 &=\gate(c,s,L_s,\{t\},\{q\},\switch(u,K_a,K_h)).
\end{align*}
Read this from the outside inward: scan for secrets; if the scanner passes, route to the AI or human reviewer; if that reviewer escalates, deny. The intended routing policy might send high-risk changes to the human, but the proof must cover either branch.

This term uses all four constructors. Its root is a gate, the gate's child is a switch, the switch has two guard children, and each guard has a \emph{syntactic} $\deny$ leaf. An execution can terminate at a guard before visiting that leaf.

\textbf{Three possible executions.}
\begin{enumerate}
\item Scanner fails: immediate denial; no reviewer is invoked.
\item Scanner passes, router selects AI, AI accepts: immediate acceptance at $a$; the $\deny$ continuation is not run.
\item Scanner passes, router selects human, human escalates: run its continuation $\deny$, so deny the proposal.
\end{enumerate}
These describe control flow. They do not yet establish safety: that also needs the contracts and their assumptions to hold.

\textbf{Now walk the static checker.} Start with $A=B=\emptyset$. The scanner's requirement $\{t\}$ is available from its own declaration $\Omega_c=\{t\}$, so it is credited. At the switch,
\[
A=\{s\},\qquad B=\{q\}.
\]
Both branches receive these same sets. At the AI reviewer, $\Omega_a=\emptyset$, but $\req(L_a)=\{q\}\subseteq\Omega_a\cup B$; its scope covers the remaining goal $r$. At the human, $q$ comes from $B$ and \ob{human\_attentive} from its local evidence; its scope also covers $r$. Both denial continuations check, so the hierarchy checks.

This verdict is conditional: the scanner's declaration $t$, the human's attention declaration, and the gate's promised goal and provided fact must be justified. The walk has not measured scanner accuracy or guaranteed that any particular reviewer will accept.
\end{example}

\begin{warning}{Three distinctions to keep visible}
\textbf{$\varphi$ versus $D$:} the scanner's goal $s$ goes into $A$; its obligation $q$ goes into $B$. A goal already covered and an assumption needed by a later contract have different jobs.

\textbf{$\Omega$ versus $D$:} the scanner declares $t$ only for its own invocation. The fact $q$ is explicitly provided to later invocations after a contract-applicable pass. Nothing automatically copies $t$ to the reviewers.

\textbf{Routing versus evidence:} selecting $K_a$ does not establish ``low risk''. A switch has no contract. If the AI reviewer needs a low-risk obligation, that obligation needs its own justification even on the branch intended for low-risk inputs.
\end{warning}

\textbf{What if the AI should escalate to a human?} Replace $K_a$ by
\[
\guard(a,L_a,\emptyset,
       \guard(h_a,L_h,\{\ob{human\_attentive}\},\deny)).
\]
If $a$ accepts or blocks, $h_a$ never runs. Only escalation reaches $h_a$, with the same $A,B$. The checker must justify acceptance at \emph{both} guards; a human continuation cannot repair an unsafe immediate acceptance by the AI.

\pauseq{In $H_{\mathrm{review}}$, remove the scanner's provided fact $q$ but leave its goal $s$. Do both reviewers still check? Would routing to the human establish that $q$ holds?}{No. The scanner still adds $s$ to $A$, but $B$ remains empty. Both reviewers require $q$, which neither declares locally, and $r$ is still open. Thus both fail the guard check. A switch changes neither $A$ nor $B$; its branch choice establishes no obligation.}

\textbf{Connection to structural induction.} The toy payment language had three agent cases and one sequencing case. This calculus instead has four proof cases: $\deny$ cannot accept; a credited passing gate justifies updated goals and facts for its continuation; a guard must be safe if it accepts and have a safe continuation if it escalates; a switch needs both branches checked because either may run. These are exactly the cases in the soundness proof of Chapter~\ref{ch:check}.

\subsection{DevAgent as a term}

Here is the running example written as a tree. Each node of the tree is one constructor.

\begin{figure}[h]
\centering
\begin{tikzpicture}[font=\scriptsize,
  gt/.style={draw=tintblue,rounded corners=2pt,fill=tintblue!8,align=center,minimum height=0.7cm,minimum width=1.7cm},
  gd/.style={draw=tintorange,rounded corners=2pt,fill=tintorange!12,align=center,minimum height=0.7cm,minimum width=1.7cm},
  sw/.style={draw=tintpurple,rounded corners=2pt,fill=tintpurple!10,align=center,minimum height=0.7cm,minimum width=1.7cm},
  dn/.style={draw=gray,rounded corners=2pt,fill=gray!15,align=center,minimum height=0.55cm,minimum width=1.0cm},
  arr/.style={-{Stealth[length=1.8mm]}}]
\node[gt] (g1) at (0,0) {gate $g_1$\\establishes $\varphi_1$};
\node[gt] (g2) at (0,-1.1) {gate $g_2$\\establishes $\varphi_2$};
\node[gt] (g3) at (0,-2.2) {gate $g_3$\\establishes $\varphi_{\mathrm{env}}$};
\node[sw] (s) at (0,-3.3) {switch $s$\\``tainted?''};
\node[gt] (v2) at (-3.3,-4.6) {veto gate $v_2$\\scorer, $\varphi=\bot$};
\node[gd] (n4) at (-3.3,-5.8) {guard $n_4$ (human)\\scope $\{\varphi_3,\varphi_4\}$};
\node[dn] (d1) at (-3.3,-6.9) {deny};
\node[gt] (c) at (3.3,-4.6) {gate $c$ (det.)\\establishes $\varphi_3$\\provides \ob{safe\_on\_this\_text}};
\node[gd] (n3) at (3.3,-5.9) {guard $n_3$ (judge)\\scope $\{\varphi_4\}$};
\node[dn] (d2) at (3.3,-7.0) {deny};
\draw[arr] (g1)--(g2); \draw[arr] (g2)--(g3); \draw[arr] (g3)--(s);
\draw[arr] (s)--node[above left,font=\tiny]{yes (tainted)}(v2); \draw[arr] (s)--node[above right,font=\tiny]{no}(c);
\draw[arr] (v2)--(n4); \draw[arr] (n4)--node[left,font=\tiny]{escalate}(d1);
\draw[arr] (c)--(n3); \draw[arr] (n3)--node[right,font=\tiny]{escalate}(d2);
\end{tikzpicture}
\caption{DevAgent as a term. Tainted input (text from outside) goes to a scorer and then a person. Untainted input is checked by a deterministic provenance gate $c$ (which \emph{provides a fact}) and then reaches the raw-text LLM judge. Denial edges for failed gates are omitted.}
\label{fig:devagent-term}
\end{figure}

Why this design? Recall the lead-lined suitcase. The judge $n_3$ requires \ob{safe\_on\_this\_text}. Gate $c$ \emph{verifies by provenance} that no untrusted text is present in the input, a purely mechanical property (a label store), and only if it passes does it \emph{provide} \ob{safe\_on\_this\_text} to later stages. The judge never sees tainted text. Tainted text goes to a human.

\pauseq{In Figure~\ref{fig:devagent-term}, which nodes are \emph{routes}? Which nodes are deterministic gates? Which node is a veto stage? How many different paths lead from $g_1$ to an accepting guard?}{Routes (guards): $n_3$ and $n_4$. Deterministic gates: $g_1,g_2,g_3,c$. Veto stage: $v_2$. Two paths reach an accepting guard: $g_1g_2g_3\,s\,v_2\,n_4$ and $g_1g_2g_3\,s\,c\,n_3$.}

\subsection{What ``running'' a hierarchy means: models}

\textbf{Why we need models.} We do not know how the stages behave. So we define a \textbf{model} to be a \emph{complete description of one possible behavior for one proposal}. Nothing is excluded unless a contract excludes it.

\begin{defbox}{model}
A \textbf{model} $M=(P,V,W,E,S)$ for one proposal consists of
\begin{itemize}
\item $P(n)\in\{\text{true},\text{false}\}$: does gate $n$ pass the proposal?
\item $V(n)\in\{\acc,\blk,\esc\}$: guard $n$'s verdict.
\item $W(n)\in\{\text{true},\text{false}\}$: which branch does switch $n$ take?
\item $E(n)$: the set of obligations that \emph{actually hold} at invocation $n$.
\item $S\subseteq\Goals$: the goals the proposal \emph{actually satisfies}.
\end{itemize}
\end{defbox}

All components are \textbf{free}: any combination of values is a possible model unless a contract excludes it (we define that in a moment). An \textbf{attack} is then simply a model: typically one in which some obligation fails at some stage, or in which a goal is not in $S$.

\begin{example}{models for the toy hierarchy}
Throughout the next chapters we use a very small hierarchy $H_{\mathrm{toy}}$ to try ideas. There are two goals $\Goals=\{a,b\}$, one obligation $o$, and one layer $L$ with $\req(L)=\{o\}$, $\scope(L)=\{b\}$.
\[
H_{\mathrm{toy}}=\gate(g,a,\epsilon,\emptyset,\emptyset,\ \guard(j,L,\Omega_j,\deny))
\]
A deterministic gate $g$ establishes $a$; then a guard $j$ (say, an LLM judge) with contract $L$ covers $b$. The evidence $\Omega_j$ is a parameter we vary.

A model: $P(g)=\text{true}$, $V(j)=\acc$, $E(j)=\{o\}$, $S=\{a,b\}$. A different model: $P(g)=\text{true}$, $V(j)=\acc$, $E(j)=\emptyset$, $S=\{a\}$. In the second, the obligation $o$ fails at $j$, and the proposal does not satisfy $b$.
\end{example}

\subsection{The decision semantics}
\label{sec:reading-semantics}

\begin{toolbox}{reading the double brackets}
The notation $\sem{H}{M}$ is called the \textbf{semantic interpretation} of $H$ under $M$. The double brackets mean ``interpret this hierarchy''; the subscript says which model supplies the behavior. They are not a multiplication or a new kind of hierarchy.

Think of this as a function $\mathrm{decision}(H,M)$: $H$ supplies the wiring, and $M$ supplies whether each gate passes, each guard's verdict, and each switch's branch. The result is $\denied$ or $\accepted(r)$, naming the guard that accepted. For example, $\accepted(j)$ says \emph{who accepted}, not \emph{which goals hold}.
\end{toolbox}

Given a model, what is the decision on the proposal? This is defined by recursion on the term (the technique of Chapter~\ref{ch:toolkit}), one equation per constructor. We write $\sem{H}{M}$ for the outcome, which is either $\denied$ or $\accepted(r)$, where $r$ is the guard that accepted.

\begin{align*}
\sem{\deny}{M}&=\denied\\
\sem{\gate(n,\varphi,\ell,\Omega,D,k)}{M}&=\sem{k}{M}\ \text{ if }P(n),\ \text{else }\denied\\
\sem{\guard(n,\ell,\Omega,k)}{M}&=\begin{cases}\accepted(n)&\text{if }V(n)=\acc\\\sem{k}{M}&\text{if }V(n)=\esc\\\denied&\text{if }V(n)=\blk\end{cases}\\
\sem{\switch(n,a,b)}{M}&=\sem{a}{M}\ \text{ if }W(n),\ \text{else }\sem{b}{M}
\end{align*}

This is just a program: follow the structure and do what each node says.

Only $P,V,W$ affect this decision. The actual obligations $E(n)$ and actual goals $S$ are used to assess contracts and safety; the decision equations do not consult them. In particular, changing whether a proposal is safe need not change the decision of a guard that accepts it.

\begin{example}{evaluating $H_{\mathrm{toy}}$}
With the first model above: $P(g)$ true so we continue to the guard; $V(j)=\acc$, so the result is $\accepted(j)$. With a model where $P(g)=\text{false}$, the result is $\denied$ regardless of what $j$ would say.

Expand the recursive calculation for a passing gate and accepting judge:
\[
\sem{H_{\mathrm{toy}}}{M}
 =\sem{\guard(j,L,\Omega_j,\deny)}{M}
 =\accepted(j).
\]
The first equality uses $P(g)=\mathrm{true}$; the second uses $V(j)=\acc$. If instead $V(j)=\esc$, the second expression becomes $\sem{\deny}{M}=\denied$. If $V(j)=\blk$, it is immediately $\denied$. No step asks whether $o\in E(j)$ or $b\in S$.
\end{example}

\subsection{Contract satisfaction: which models are allowed}

The decision semantics says what \emph{would happen}. But models in which a guard accepts everything are possible, and for those the question ``is it safe?'' would be hopeless. The contracts restrict which models count. We define the relation
\[
M\models_F H
\]
read: ``this behavior respects every contract on the invocations the proposal reaches, given the facts $F$ so far''.
The symbol $\models$ is read ``satisfies''; $M\not\models_F H$ means the statement is false. Unlike the decision semantics, this relation returns no verdict and does not choose the behavior. It asks whether the behavior already specified by $M$ meets the contract rules, including realization.

Here $F$ is the set of facts provided so far on this modeled run. The initial safety definition uses $F=\emptyset$; safety from $(A,B)$ uses $F=B$, and the generalized proof allows any $F\supseteq B$. The relation is again defined by one equation per constructor:

\begin{align*}
M\models_F&\ \deny\quad\text{always}\\
M\models_F&\ \gate(n,\varphi,\ell,\Omega,D,k)\ \text{ iff, when }P(n):\ \Real(n)\ \wedge\\
&\quad\text{if }\req(\ell)\subseteq E(n):\ \varphi\in S\ \wedge\ M\models_{F\cup D}k\\
&\quad\text{otherwise: }\ M\models_F k\\
M\models_F&\ \guard(n,\ell,\Omega,k)\ \text{ iff }\ \Real(n)\ \wedge\\
&\quad(V(n)=\acc\ \wedge\ \req(\ell)\subseteq E(n)\Rightarrow\scope(\ell)\subseteq S)\\
&\quad\wedge\ (V(n)=\esc\Rightarrow M\models_F k)\\
M\models_F&\ \switch(n,a,b)\ \text{ iff }\ M\models_F a\text{ if }W(n),\ \text{else }M\models_F b
\end{align*}
where $\Real(n)\triangleq\Omega\cup F\subseteq E(n)$ is called \textbf{realization}.

This looks dense, so let us read it slowly. There are three ideas, matching the three modelling choices from Chapter~\ref{ch:toolkit}.

\begin{enumerate}
\item \textbf{Realization.} At each reached guard and each passing gate, the \emph{declared} evidence $\Omega$ and the \emph{provided} facts $F$ must be \emph{true}, i.e.\ belong to $E(n)$. A false declaration there puts the model outside the guarantee: it is not a ``contract-respecting behavior''. A failed gate already denies and its satisfaction rule imposes no further condition. Declared evidence that the checker cannot verify ends up on a list of assumptions: you are asserting that it holds.
\item \textbf{Conditionality.} Look at the guard line: ``$V(n)=\acc\wedge\req(\ell)\subseteq E(n)\Rightarrow\scope(\ell)\subseteq S$''. The promise (scope is satisfied) is only required when \emph{all} the obligations really hold at this invocation. Otherwise nothing is demanded: the guard can accept something unsafe and still be ``respecting its contract''.
\item \textbf{Locality.} $E$ is indexed by the node $n$. What holds at the scorer says nothing about what holds at the judge. And a fact in $F$ reaches later stages \emph{only} through a gate whose contract applied. Note the guard line: for $V(n)=\esc$ the rest is checked under the same $F$, not a larger one. A guard adds nothing to $F$.
\end{enumerate}

\begin{example}{contract satisfaction for $H_{\mathrm{toy}}$}
Let $\Omega_j=\emptyset$ (the specification does not promise that $o$ holds at $j$). Take the second model above: $P(g)$, $V(j)=\acc$, $E(j)=\emptyset$, $S=\{a\}$.
\begin{itemize}
\item Gate $g$ ($\epsilon$ layer): $\req=\emptyset\subseteq E(g)$, so the contract applies; we need $a\in S$ and it is. \checkmark
\item Guard $j$: realization $\Omega_j\cup F=\emptyset\subseteq E(j)$ \checkmark. Acceptance: is $\req(L)=\{o\}\subseteq E(j)=\emptyset$? No. So the condition is false and the implication is vacuously true. \checkmark
\end{itemize}
So $M\models_\emptyset H_{\mathrm{toy}}$: this \emph{is} a contract-respecting behavior. Yet $\sem{H_{\mathrm{toy}}}{M}=\accepted(j)$ and the proposal violates $b$. This is the lead-lined suitcase in miniature. If we instead declared $\Omega_j=\{o\}$, then realization would demand $o\in E(j)$, and this model would be excluded.
\end{example}

\subsubsection{Same decision, different contract judgments}
\label{sec:contracts-versus-safety}

Give $o$ a concrete reading: ``the judge's input satisfies the sanitization policy''. The judge promises goal $b$ only when $o$ holds and it accepts. In all four rows below, use the toy hierarchy, let the deterministic gate pass, and let the judge accept. Set $E(g)=\emptyset$; every listed $S$ contains the gate's promised goal $a$. No gate provides facts, so $F=\emptyset$ at $j$.

\begin{center}
\renewcommand{\arraystretch}{1.25}
\begin{tabularx}{\textwidth}{@{}llllX@{}}
\toprule
$\Omega_j$ & $E(j)$ & $S$ & \textbf{Contracts?} & \textbf{Reason}\\
\midrule
$\emptyset$ & $\{o\}$ & $\{a,b\}$ & Respected & The assumption holds and the judge's promised goal $b$ holds.\\
$\emptyset$ & $\{o\}$ & $\{a\}$ & Violated & The assumption holds, but the accepting judge fails its promise.\\
$\emptyset$ & $\emptyset$ & $\{a\}$ & Respected & Realization holds; the missing, undeclared assumption makes the judge's promise inapplicable.\\
$\{o\}$ & $\emptyset$ & $\{a\}$ & Violated & Realization fails: the local declaration says $o$ holds, but it does not.\\
\bottomrule
\end{tabularx}
\end{center}

Every row has the same decision $\sem{H_{\mathrm{toy}}}{M}=\accepted(j)$. The third row is the crucial unsafe, contract-respecting behavior. The fourth row is excluded because its declaration is false, even though the conditional promise alone would be vacuously true. \textbf{A missing assumption and a false evidence declaration are different failures.}

\begin{warning}{Respecting a contract does not mean its assumption held}
The relation $M\models_F H$ checks both realization and the conditional promises. An obligation absent from $E(n)$ can make a promise inapplicable without breaking the contract, provided realization does not require that obligation. A fact declared in $\Omega$ or carried in $F$ must hold wherever realization is required. This is why ``the contracts are respected'' and ``the accepted proposal is safe'' are separate statements.
\end{warning}

\pauseq{In the third row, why do acceptance and contract satisfaction both hold even though $b$ fails? What changes if $o$ is added to $\Omega_j$?}{Acceptance follows from $P(g)=\mathrm{true}$ and $V(j)=\acc$, independently of $E$ and $S$. Realization is $\emptyset\subseteq E(j)$, and the judge's conditional promise has a false assumption, so contracts are respected. Adding $o$ to $\Omega_j$ makes realization fail in this same model; it does not sanitize the input or change the actual verdict.}

\begin{warning}{Do not mix up the three things}
There are three different ``kinds'' of an obligation $o$, and confusing them is the number-one source of errors when reading this paper:
\begin{center}
\renewcommand{\arraystretch}{1.25}
\begin{tabular}{@{}lll@{}}
\toprule
\textbf{Kind} & \textbf{Where it lives} & \textbf{Who knows it}\\
\midrule
\emph{Declared} ($o\in\Omega$) & the specification (syntax) & the checker\\
\emph{True} ($o\in E(n)$) & the model (the world) & nobody; this is what must be argued\\
\emph{Credited} ($o\in B$) & the analysis & the checker\\
\bottomrule
\end{tabular}
\end{center}
Realization is the bridge between the first two: declaring something makes it a \emph{premise}. Soundness (Chapter~\ref{ch:check}) is the bridge between the last two and the world: whatever the checker credits really is true in every contract-respecting model.
\end{warning}

\subsection{Safety: the definition everything is about}

\begin{defbox}{safety}
$H$ is \textbf{safe from} $(A,B)$ if for every model $M$ with $A\subseteq S$ and $M\models_B H$: if $\sem{H}{M}=\accepted(r)$ for some $r$, then $\Goals\subseteq S$. $H$ is \textbf{safe} if it is safe from $(\emptyset,\emptyset)$.
\end{defbox}

In words: \emph{take any behavior that respects all contracts. If the proposal is accepted, then all the goals really hold.} The extra $(A,B)$ say ``suppose we already know the goals $A$ hold and the facts $B$ have been provided''; they let us talk about the \emph{tail} of a hierarchy after some gates have passed, which is essential for induction.

For the initial state, unpack the definition as
\[
\forall M\,\forall r:\quad
\bigl(M\models_\emptyset H\ \wedge\ \sem{H}{M}=\accepted(r)\bigr)
\ \Longrightarrow\ \Goals\subseteq S.
\]
Read left to right: choose any model and guard; restrict attention to contract-respecting behaviors; if that behavior accepts at the guard, require every goal to hold. The set $S$ belongs to the chosen model $M$. One safe accepted model does not prove this universal claim. The third row above \emph{refutes} safety of the toy hierarchy with $\Omega_j=\emptyset$, because it supplies a single contract-respecting accepted model with $b\notin S$.

\begin{center}
\begin{tabularx}{\textwidth}{@{}p{3.4cm}XX@{}}
\toprule
\textbf{Expression} & \textbf{Question} & \textbf{Input and result}\\
\midrule
$\sem{H}{M}$ & What decision happens? & A hierarchy and one model; returns $\denied$ or $\accepted(r)$.\\
$M\models_F H$ & Does this behavior respect the contract rules? & One model, a hierarchy and starting facts; true or false.\\
$\chk(H,A,B)$ & Do the declared contracts suffice for safe composition? & A specification and credited sets, without choosing a model; true or false.\\
\bottomrule
\end{tabularx}
\end{center}
Soundness connects the third row to the first two: if the checker succeeds, every contract-respecting behavior that accepts satisfies the goals. It does not assert that the system will accept a proposal.

Observe four things about this definition.

\begin{enumerate}
\item It is about \emph{one proposal at a time}, and about \emph{all behaviors} at once. No probabilities yet.
\item It is \emph{conditional} on the contracts. An attacker who makes a contract false is outside it.
\item It says nothing about usefulness: $\deny$ is safe (nothing is ever accepted). A hierarchy that rejects everything is vacuously safe, and the calculus would approve it.
\item It is about \emph{acceptance}: if the proposal is denied, nothing is claimed. This matches our concern: the damage is done by accepted actions.
\end{enumerate}

\begin{table}[h]
\caption{What remains against each attacker class (paper, Table~I).}
\label{tab:attackers-tut}
\centering
\footnotesize
\renewcommand{\arraystretch}{1.25}
\begin{tabularx}{\textwidth}{@{}p{3.5cm}XX@{}}
\toprule
\textbf{Attacker} & \textbf{Goals set by deterministic gates} & \textbf{Goals resting on learned stages}\\
\midrule
None; drift only & hold exactly & certificate degrades by the drift $\eta$\\
Fixed attacks, present in calibration & hold exactly & certificate holds for that mix\\
Adaptive & hold, given complete mediation & no guarantee; the contract may be false\\
\bottomrule
\end{tabularx}
\end{table}

The ``drift'' $\eta$ is the \emph{total-variation distance} between the data distribution used for calibration and the one in deployment (how much the world changed); we meet it in Chapter~\ref{ch:risk}. ``Complete mediation'' means every effect of the agent really passes through the gates; if the agent has a raw shell the gates do not see, nothing here applies.

\subsection{Quick reference: all the sets}

\begin{table}[h]
\caption{Sets used in the calculus (paper, Table~II).}
\centering
\footnotesize
\renewcommand{\arraystretch}{1.2}
\begin{tabularx}{\textwidth}{@{}lX@{}}
\toprule
\textbf{Set} & \textbf{Meaning}\\
\midrule
$\Omega$ & obligations \emph{declared} at this stage (syntax)\\
$D$ & facts a gate \emph{provides} (syntax)\\
$E(n)$ & obligations \emph{actually true} at invocation $n$ (model)\\
$S$ & goals the proposal \emph{actually satisfies} (model)\\
$F$ & facts \emph{actually provided} so far on the run (semantics)\\
$A,B$ & goals and facts the \emph{checker has credited} so far (analysis)\\
\bottomrule
\end{tabularx}
\end{table}

The analysis tracks declarations; soundness shows that $A\subseteq S$ and $B\subseteq F$ are preserved along a run. That sentence is the heart of Chapter~\ref{ch:check}.

\subsubsection{Follow the three records through a complete example}
\label{sec:accounting-example}

\begin{example}{a sanitization validator, a content filter, and a judge}
Use goals $\Goals=\{c,m,r\}$: the proposal's judge-facing view satisfies the sanitization policy ($c$), it satisfies the content policy ($m$), and it meets the final review goal ($r$). Let $p=\ob{sanitized\_view}$ be an \emph{obligation} needed by later contracts. Its intended meaning is related to $c$, but the calculus does not infer $p$ from the goal name $c$: the validator must explicitly provide it.

The validator checks an already prepared view of the same proposal; it does not rewrite the proposal during this run. Define
\[
\begin{array}{c|cc}
 & \req & \scope\\\hline
L_m & \{p\} & \emptyset\\
L_j & \{p\} & \{r\}
\end{array}
\]
and the hierarchy
\begin{align*}
K_j&=\guard(j,L_j,\emptyset,\deny),\\
K_m&=\gate(f,m,L_m,\emptyset,\emptyset,K_j),\\
H_{\mathrm{san}}&=\gate(v,c,\epsilon,\emptyset,\{p\},K_m).
\end{align*}
The deterministic validator $v$ establishes $c$ and provides $p$ if it passes. The content filter $f$ requires $p$, establishes $m$ if it passes with its contract applicable, and provides no new obligations. The judge $j$ requires $p$ and covers $r$ on acceptance.

For comparison with the static walk, take a contract-respecting run in which both gates pass and the judge accepts:
\begin{center}
\renewcommand{\arraystretch}{1.25}
\begin{tabularx}{\linewidth}{@{}lcccX@{}}
\toprule
\textbf{Point} & $A$ & $B$ & $F$ & \textbf{What justifies the next step?}\\
\midrule
Start & $\emptyset$ & $\emptyset$ & $\emptyset$ & The validator has no required obligations.\\
After $v$ & $\{c\}$ & $\{p\}$ & $\{p\}$ & $p$ is available to the filter's contract.\\
After $f$ & $\{c,m\}$ & $\{p\}$ & $\{p\}$ & The filter adds a goal, not a new provided fact.\\
At $j$ & $\{c,m\}$ & $\{p\}$ & $\{p\}$ & $p$ makes the judge's contract usable; its scope covers $r$.\\
\bottomrule
\end{tabularx}
\end{center}
The checker credits goals in $A$ and obligations in $B$. The semantic relation carries the corresponding provided obligations in $F$. On acceptance the judge's contract gives $r\in S$; the gates have already given $c,m\in S$, so $\Goals\subseteq S$. The judge's scope does not get added to $A$ for a continuation: acceptance ends the run, and escalation adds nothing.

Passing the validator supports these conclusions only under its promised contract, including the claim that $p$ holds at later invocations. Merely naming an obligation \ob{sanitized\_view} does not implement sanitization or verify that claim.
\end{example}

\begin{samepage}
\textbf{The bridge from records to truth.} At a reached guard or a passing gate $n$, the soundness proof uses
\[
\underbrace{B\subseteq F}_{\text{credited facts are provided}}
\quad\text{and}\quad
\underbrace{\Omega_n\cup F\subseteq E(n)}_{\text{realization}}
\quad\Longrightarrow\quad
\Omega_n\cup B\subseteq E(n).
\]
\end{samepage}
Therefore a requirement available in $\Omega_n\cup B$ really holds at that invocation in a contract-respecting model. Separately, $A\subseteq S$ connects credited goals to goals actually satisfied. After a credited passing gate, its contract gives $\varphi\in S$, so $A\cup\{\varphi\}\subseteq S$ remains true; adding the same $D$ to $B$ and $F$ preserves $B\cup D\subseteq F\cup D$.

\begin{warning}{The semantic record is not a runtime measurement}
$F$ is the parameter carried by $M\models_F H$; it is not an independently observed log or an extra argument to the static checker. $E(n)$ describes the obligations actually holding at a particular invocation; $F$ contains the obligations earlier applicable gates provided, which realization requires to hold there. $E(n)$ can contain more than $F$.

A passing gate adds $D$ to $F$ only when $\req(\ell)\subseteq E(n)$. The static checker adds $D$ to $B$ only when $\req(\ell)\subseteq\Omega\cup B$. These tests use different information. Thus $B$ can be smaller than $F$, as the soundness example in Chapter~\ref{ch:check} shows. If the actual contract is inapplicable, passing alone supplies no new semantic facts.
\end{warning}

\pauseq{In $H_{\mathrm{san}}$, remove $p$ from the validator's provided set while keeping its goal $c$. May the checker still use the content filter's and judge's contracts?}{No. After $v$, $A=\{c\}$ but $B=\emptyset$. Neither later stage declares $p$, so neither contract can be credited from $\Omega\cup B$. The checker does not convert goal $c$ into obligation $p$. The content filter contributes no credited goal, and the judge cannot cover the open goals. A particular model might still have $p\in E(n)$, but the static checker must justify all contract-respecting behaviors, not just that model.}

\subsection{The concrete language: \Hsl{}}

The calculus is a theory object. The tool works on \Hsl{}, a JSON file format that describes the same thing. Instead of a term it describes a \emph{graph}: nodes may share successors (two branches that rejoin), and the term is the \emph{unfolding} along paths. Chapter~\ref{ch:graphs} explains why this matters for speed. The excerpt below is from the DevAgent file in the repository.

\begin{lstlisting}[language={},basicstyle=\ttfamily\scriptsize]
"layers": [{"name": "L3", "scope": ["phi4"], "view": "raw",
  "requires": ["miss_rate_bounded", "config_match"]}, ...],
"nodes": [ ...,
 {"type": "switch", "id": "s", "condition": "input_tainted",
  "then": "v2", "otherwise": "c"},
 {"type": "gate", "id": "c", "rule": "no_untrusted_input",
  "establishes": "phi3", "next": "n3",
  "provides": ["safe_on_this_text"]},
 {"type": "guard", "id": "n3", "layer": "L3", "next": null,
  "evidence": ["miss_rate_bounded", "config_match"]}, ...]
\end{lstlisting}

You can match every field to the calculus: \code{requires} is $\req$, \code{scope} is $\scope$, \code{evidence} is $\Omega$, \code{provides} is $D$, \code{establishes} is $\varphi$, \code{next:\ null} is $\deny$.

\subsubsection*{Three kinds of obligations}

The language labels each obligation with a \emph{kind}, which records how it gets established (it does \emph{not} change the safety check; it changes what is left to argue):

\begin{center}
\renewcommand{\arraystretch}{1.25}
\begin{tabularx}{\textwidth}{@{}lXX@{}}
\toprule
\textbf{Kind} & \textbf{Meaning} & \textbf{Example}\\
\midrule
\textbf{checked} & has a runtime predicate; failing it denies & \ob{config\_match}: the deployed judge uses the configuration that was certified\\
\textbf{construction} & true of the code, e.g.\ by how an input is built & \ob{human\_sees\_provenance}: the interface always shows where the text came from\\
\textbf{assumed} & the deployer's responsibility to justify & \ob{miss\_rate\_bounded}, \ob{human\_attentive}, \ob{safe\_on\_this\_text}\\
\bottomrule
\end{tabularx}
\end{center}

Obligations may also be tied to a \emph{view} (the input a layer sees, say ``raw'' vs.\ ``structured''); evidence tied to another view than the stage's layer is dropped, with a warning, before the analysis, because it describes a different invocation (locality again).

\begin{recap}
\begin{itemize}
\item Four building blocks: $\deny$, gate (can only stop; can establish a goal and provide facts), guard (can accept, block, or escalate; provides nothing), switch (routes; teaches nothing).
\item A \emph{model} is a free description of one behavior. An attack is a model.
\item $\sem{H}{M}$ is what happens; $M\models_F H$ says all contracts on the reached stages are respected, with realization, conditionality and locality.
\item \emph{Safe}: every contract-respecting model that is accepted satisfies all goals.
\item Distinguish \emph{declared}, \emph{true} and \emph{credited}.
\end{itemize}
\end{recap}
